\SourcePage{419}{422}
then be interpreted visually as follows. Equal potential (and total) energies also imply equal momenta. Thus, in place of (90.1), one may write
\begin{equation}r_1=r_2,\qquad p_1=p_2,\tag{90.2}\end{equation}
where $p$ is the momentum of the relative radial motion of the nuclei, while the subscripts 1 and 2 refer to the two electronic states. Thus, at the instant of transition, the separation and momentum of the nuclei may be said to remain unchanged (the so-called \emph{Franck--Condon principle}). Physically, this follows because electronic velocities are large in comparison with nuclear velocities, and ``during the electronic transition'' the nuclei have no time to change their positions and velocities appreciably.
